有一类线上问题的症状高度一致：同一套阈值、同一份配置，在某些时段安静得像坏了，在另一些时段一天报警几十次。查下去会发现被监控的量并没有跑偏，均值稳稳的，只是散布度换了档。收益率、请求量、端到端延迟、模型错误率，都会出现这个形态——大波动扎堆，平静也扎堆。

用一阶的工具去看，什么也看不见。序列的自相关图整排落在显著性带内，任何以``可预测''为名的检查都会给出否定的结论。但把每个值平方一遍再画自相关图，柱子会整排站到带外，而且衰减得很慢。方向不可预测，幅度可预测，这两句话在同一份数据上同时成立。这就是条件异方差留下的指纹：独立同分布的假设在二阶矩上已经破了，而一阶看不出来。

拟合一个 GARCH(1,1)，典型的结果是 \(\hat\alpha\approx0.10\)、\(\hat\beta\approx0.85\)。两个参数各自不起眼，但它们的和 \(\hat\alpha+\hat\beta\approx0.95\) 才是要看的数：它决定一次冲击的半衰期，这里是十三四期。也就是说，一次事故过后接近三周内的风险敞口都高于常态，而按长期均值做的容量与阈值设定，会在这三周里持续偏低。

本单元的三篇按同一条链条推进。第一篇建立诊断动作：拟合任何时间序列模型之后，除了看残差的自相关，还要看残差平方的自相关，并给出 ARCH-LM 这样的正式检验。第二篇建立模型：ARCH 让条件方差依赖过去的冲击平方，GARCH 再让它依赖自己的过去，从而用两三个参数表达一条平滑衰减的记忆曲线——和 ARMA 用少数参数表达长记忆是同一个手法；平稳性条件 \(\alpha+\beta<1\) 在这里是判断模型是否可信的第一道关，等于一就是 IGARCH，冲击永不消失。第三篇处理落地：条件似然怎么写、标准化残差怎么诊断、杠杆效应与厚尾新息各自补上哪块短板。

有一条贯穿全篇的界线值得先立起来：GARCH 建模的是条件方差，不是收益率的方向。它不告诉你明天涨还是跌，只告诉你明天的分布有多宽。这个区分决定了它的正确用法在风险度量、区间宽度、动态阈值和容量余量这一侧，而不在方向性决策那一侧。

\subsection{怎么读这一章}

从第一篇的经验事实进入最省力：先确认自己手上的数据是不是这个形态，再决定要不要往下走。若已经确认，第二篇的重点只有一个数——\(\alpha+\beta\)，把它的含义吃透，模型的行为基本就在预料之中了。第三篇偏工程，可以按需查阅：估计不收敛看第一小节，残差还有结构看第二小节，尾部算不准看厚尾那一小节。

\subsection{本章篇目}

以下篇目按概念依赖排列；若从具体问题进入，也可以先打开最接近当前困惑的一篇，再沿篇内链接回补。

\begin{itemize}
\tightlist
\item \hyperref[sec:11-Mathematical-Statistics/10-GARCH-and-Volatility/01]{``条件异方差的统计特征''}；
\item \hyperref[sec:11-Mathematical-Statistics/10-GARCH-and-Volatility/02]{``ARCH 与 GARCH 的随机过程结构''}；
\item \hyperref[sec:11-Mathematical-Statistics/10-GARCH-and-Volatility/03]{``GARCH 的估计、诊断与扩展''}。
\end{itemize}

读完这三篇，判断标准很简单：拿到一段序列，你能说清它的波动是否有记忆、记忆有多长、以及这段记忆该怎样改写你手上的阈值和余量。
